% DOCUMENT CLASS SETUP
\documentclass[12pt, a4paper]{article}

% DOCUMENT PACKAGES
\usepackage[a4paper, margin=1in]{geometry}
\usepackage{titling}
\usepackage{enumitem}
\usepackage{amsmath}
\usepackage{amssymb}
\usepackage{mathtools}
\usepackage{graphicx}
\usepackage{xcolor}
\usepackage{listings}
\usepackage{hyperref}
\usepackage{csquotes}
\usepackage{tocloft}

% FONT AND ENCODING SETUP
\usepackage{fontspec}
\usepackage{polyglossia}
\usepackage[Thai]{ucharclasses}
\setdefaultlanguage{english}
\setotherlanguage{thai}
\setTransitionsFor{Thai}{\thaifont}{\normalfont}
\newfontfamily{\thaifont}{THSarabun.ttf}[
    Path=./,
    BoldFont=THSarabunBold.ttf,
    ItalicFont=THSarabunItalic.ttf,
    BoldItalicFont=THSarabunBoldItalic.ttf,
    Scale=1.23
]
\XeTeXlinebreaklocale "th"
\XeTeXlinebreakskip = 0pt plus 1pt

% PROFESSIONAL CODE BLOCK
\usepackage{tcolorbox}
\tcbuselibrary{minted, skins, breakable}

% DEFINE GITHUB DARK THEME COLORS %
\definecolor{ghback}{HTML}{0d1117}    % Dark background
\definecolor{ghtext}{HTML}{c9d1d9}    % Light gray text
\definecolor{ghborder}{HTML}{30363d}  % GitHub Border color

% CODE BLOCK CONFIGURATION (MINTED) %
\renewcommand{\theFancyVerbLine}{\textcolor{white!60!gray}{\scriptsize\arabic{FancyVerbLine}}}
\newtcblisting{codeblock}[2][]{
    enhanced,
    breakable,
    colback=ghback,
    colframe=ghborder,
    coltitle=ghtext,
    colupper=ghtext,
    fonttitle=\bfseries\small,
    arc=3mm,
    boxrule=0.5mm,
    % PADDING SETTINGS
    left=25pt,
    right=10pt,
    top=5pt,
    bottom=5pt,
    % MINTED SETTINGS
    listing engine=minted,
    listing only,
    minted language=python,
    minted options={
        linenos,
        numbersep=5pt,
        fontsize=\footnotesize,
        breaklines=true,
        breakanywhere=true,
        autogobble,
        style=monokai,
    },
    title={#2},
    #1
}

% DOCUMENT TITLE SETUP %
\title {
	\textbf{2110446 DATA SCIENCE AND DATA ENGINEERING} \\
    \vspace{1em}
	\Large{\textbf{Homework 9: Spatial Anaysis}} \\
    \vspace{0.5em}
    \large{Computer Engineering, Chulalongkorn University}
}
\author{Worralop Srichainont}
\date{\today}
\setlength{\droptitle}{15em}

% DISABLE PAGE NUMBER IN TABLE OF CONTENTS
\tocloftpagestyle{empty}

% DOCUMENT CONTENTS
\begin{document}

    \maketitle
    \thispagestyle{empty}
    \clearpage

    \tableofcontents
    \thispagestyle{empty}
    \clearpage

    \pagenumbering{arabic}

    \section{Source Code}

        The modified \texttt{streamlit\_airbnb.py} can be accessed via this link
        \texttt{\url{https://drive.google.com/file/d/1VbKk_hfcBod0N9phr95emXl3ewMmOHaf/view?usp=sharing}}
        or an attached file in the MyCourseVille submission.

        \vspace{1em}

        Note that the Mapbox is now requires the API key to be able to access the map visualization.
        Make sure to create a \texttt{/.streamlit/secrets.toml} file
        to make the page able to show the map beneath the scatter or heatmap plot.

        \begin{codeblock}{API Key Configuration}
            MAP_BOX_API_KEY = "[INSERT YOUR API TOKEN HERE]"
        \end{codeblock}

        You can create your own API token
        by create a free account at the website \texttt{\url{https://www.mapbox.com}}

        \vspace{1em}

        To utilize the API key, please insert the \texttt{api\_keys} parameter on the PyDeck function.
        If you don't configure the API keys and passing the API key to the PyDeck parameter,
        the map layer won't appear on the screen.

        \begin{codeblock}{Utilize the API Keys}
            st.pydeck_chart(
                pdk.Deck(
                    api_keys={"mapbox": st.secrets.get("MAP_BOX_API_KEY")},
                    layers=[heatmap_layer],
                    initial_view_state=pdk.ViewState(
                        latitude=filtered_data["latitude"].mean(),
                        longitude=filtered_data["longitude"].mean(),
                        zoom=11,
                        pitch=0,
                    ),
                    map_style=MAP_STYLES[map_style],
                ),
                height=600,
            )
        \end{codeblock}

    \pagebreak

    \section{Map Visualization}

        In the Streamlit page, there are total 4 map visualizations which are.

        \begin{figure}[H]
            \centering
            \includegraphics[width=\textwidth]{image_01}
            \caption{Accommodation Hotspot Clusters}

            \vspace{1cm}

            \includegraphics[width=\textwidth]{image_02}
            \caption{Accommodation Hotspot Heatmap}
        \end{figure}

    \pagebreak

        \begin{figure}[H]
            \centering
            \includegraphics[width=\textwidth]{image_03}
            \caption{Accommodation Hotspot Hexagon Heatmap}

            \vspace{1cm}

            \includegraphics[width=\textwidth]{image_04}
            \caption{Kernel Density Estimation (KDE) Analysis}
        \end{figure}

    \pagebreak

    \section{Answers}
        Consider only the accommodation which price is not exceed 1,000 baht. Answer the following questions

        \subsection{Problem I}
            According to the DBSCAN clustering by setting parameters \texttt{eps} to 0.001 (approximately 111 meters)
            and \texttt{min\_samples} to 3. Which district has the biggest cluster?

            \vspace{1em}

            \noindent First, set the parameters to this at the sidebar of the Streamlit page.
            \begin{figure}[h]
                \centering
                \includegraphics[width=0.3\textwidth]{image_05}
                \caption{Parameters Configuration}
            \end{figure}

            \noindent Then, go to the cluster statistics section and look at the first one.
            \begin{figure}[h]
                \centering
                \includegraphics[width=\textwidth]{image_06}
                \caption{Cluster Information}
            \end{figure}

            \noindent \underline{\textbf{Ans.}} Khlong Toei district and Vadhana district has the biggest cluster.

    \pagebreak

        \subsection{Problem II}
            According to the KDE analysis by setting parameters \texttt{bandwidth} to 0.005 (approximately 555 meters)
            as default. What is the accommodation name at the highest density point?

            \vspace{1em}

            \noindent First, set the parameters to this at the sidebar of the Streamlit page.
            \begin{figure}[h]
                \centering
                \includegraphics[width=0.3\textwidth]{image_07}
                \caption{Parameters Configuration}
            \end{figure}

            \noindent Then, go to the KDE top 10 density listings right below the map section and look at the first one.
            \begin{figure}[h]
                \centering
                \includegraphics[width=\textwidth]{image_08}
                \caption{Top 10 Density Listings}
            \end{figure}

            \noindent \underline{\textbf{Ans.}} Sukhumvit Suite: Onnut BTS, Malls \& Street Food.

\end{document}
